\section{Reasoning services and interfaces}
\label{sec:services}

\subsection{Queries beyond the class hierarchy}

An ontology editor asks a reasoner about individuals and properties as well as
classes.  KM answers these queries by reduction to classification: it extends
the input with fresh named classes whose position in the taxonomy encodes the
answer, classifies the extended ontology once, and reads the rows off the
result.  For example, the types of an individual $a$ are the named superclasses
of a fresh class equivalent to the nominal $\{a\}$, and two individuals are
equal when their nominal classes are equivalent.  The command-line options of
Table~\ref{tab:rows} select the rows.  The reductions are exact; the plugin
computes a part of some answers in Java, and the repository lists each
reduction with the part computed outside the engine.

\begin{table}[t]
\centering
\small
\caption{Result rows of \texttt{km classify}.  The class hierarchy,
unsatisfiable classes, and consistency are always reported.}
\label{tab:rows}
\begin{tabularx}{\linewidth}{@{}lX@{}}
\toprule
Option & Rows added \\
\midrule
\texttt{--realize} & the named classes each individual belongs to, and pairs of equal individuals \\
\texttt{--values} & every entailed object property assertion between named individuals, and data property values \\
\texttt{--different} & every entailed pair of different named individuals \\
\texttt{--properties} & the object property hierarchy, domains, and ranges \\
\texttt{--data-properties} & the data property hierarchy, empty data properties, domains, and disjoint data properties \\
\texttt{--top} & the classes equivalent to \texttt{owl:Thing} \\
\texttt{--export-model} & a finite model of a consistent input: elements with their classes and individuals, property edges closed under the RBox, data values, and merged individuals \\
\bottomrule
\end{tabularx}
\end{table}

\subsection{Input}

KM reads OWL~2 functional syntax directly and OWL/XML, RDF/XML, and Turtle
through a conversion to functional syntax.  In its default strict mode KM
refuses an ontology that violates the OWL~2~DL global restrictions, with a
message that names the violation.  Of the 1,920 ontologies of the ORE~2015
pool, 212 are refused for this reason, most often because anonymous individuals
break the forest condition of the structural specification.  Other reasoners
and the OWL API load many of these ontologies, so KM provides a lenient mode
that rewrites the source before reasoning, for example by replacing an
anonymous individual in a value restriction, and prints one row for each
rewrite.  A separate option treats a datatype outside the OWL~2 datatype map
as a fresh datatype with opaque literals.  The Prot\'eg\'e plugin enables both,
which matches what Prot\'eg\'e users obtain from HermiT.

\subsection{Explanations}

KM explains named-class subsumption, named-class unsatisfiability, and
ontology inconsistency in terms of source axioms.  The explanation procedure
classifies a subset of the source axioms with the automatic route, verifies
that the subset entails the query, minimises it by deletion or by delta
debugging, checks the final subset again, and uses a hitting-set tree to
search for further justifications
\cite{kalyanpur2007debugging,horridge2011justification}.  Each published
justification entails the query and loses the entailment when any one of its
axioms is deleted, which by monotonicity makes it minimal with respect to set
inclusion.  Budgets bound the number of source axioms, classification calls,
and returned justifications; when a budget is exhausted, KM discards the
unfinished justification and marks the enumeration as incomplete.

\subsection{Incremental classification}

For an ontology under edit, KM keeps a session that owns the source axioms
with stable identifiers, the normalised clauses, the selected route, the state
of the procedure, and the last answer.  An update is a transaction.  KM applies
the changed axioms to a private revision of the source, normalises the
revision, profiles and routes it, and either applies a delta to the retained
state or moves the session to another procedure.  It commits the new source,
state, and answer together, and it keeps the preceding revision unchanged when
any step fails.  The retained state depends on the procedure: the EL
procedure keeps labels, role edges, and indexes; the consequence-based engine
keeps contexts, clauses, and derivation dependencies; and the hypertableau
procedures keep completed models and their dependencies.  KM also reuses
results at the level of source modules, so that an update that touches one
part of the ontology reclassifies the affected module only.

\subsection{OWL API and Prot\'eg\'e}

A Java module implements the reasoner factory and lifecycle of the OWL API
\cite{horridge2011owlapi}, including buffering and non-buffering modes,
interruption, and disposal, and it implements the OWL Explanation API for
named subclass entailments.  The Prot\'eg\'e plugin, built for Prot\'eg\'e~5.6
and Prot\'eg\'e~6, adds KM to the Reasoner menu and to the Explain action.  The
plugin answers the methods of the reasoner interface that HermiT answers:
types, instances, equal and different individuals, the object and data
property hierarchies, domains and ranges, property values, class expression
queries as used by the DL Query tab, disjoint classes and properties,
entailment checks for the axiom types, and precomputation.  It follows
HermiT's conventions for interruption, time-outs, and inconsistent ontologies.
The plugin archive contains the \texttt{km} executable for five platforms, so
the user installs one file.  Every query runs \texttt{km} as a child process,
and the plugin terminates the process tree when the query ends, is
interrupted, or times out.
